\section{Results}

All structural values below were recomputed by the API-free benchmark from the
three frozen cases. They are exact counts for this corpus rather than estimates
of a room population. The backend tests used a deterministic response stub, and
the Unity values are editor batch-mode results. No result in this section comes
from a human-participant study or measures answer quality.

\begin{table*}[t]
  \caption{Main measured results. Counts describe the fixed three-scene corpus
  and deterministic test paths, not participants or dialogue quality.}
  \label{tab:results-summary}
  \centering
  \footnotesize
  \begin{tabular}{@{}p{.22\textwidth}p{.43\textwidth}p{.27\textwidth}@{}}
    \toprule
    Check & Observed result & Boundary \\
    \midrule
    Chains and routing
      & 93/93 assets and 186/186 target-specific pairs complete; 93/93 owner
        decisions and 459/459 route classifications agree
      & Static records; shared resolver \\
    Common injected faults
      & Both adapters detected and localized 9/9; 0/3 false positives
      & Three scripted fault types per scene \\
    Runtime admission
      & 298/298 local/direct routes admitted; 161/161 transitive/unreachable
        routes rejected before the stub
      & Deterministic response stub; restored sessions \\
    Unity and placement
      & 3/3 batch entry points, 31 contexts, and 73 actions; 4/4 placement
        workflows passed
      & Desktop Editor and scripted paths \\
    \bottomrule
  \end{tabular}
\end{table*}


\subsection{RQ1: Representation and routing conformance}

Table~\ref{tab:results-summary} gives the main counts; Appendix
Tables~\ref{tab:rq1-chains} and~\ref{tab:rq1-routing} give per-scene values.
Every one of the 93 scene assets occurred in at least one complete chain. All 186
\texttt{ScenarioStep}/\texttt{CapabilityUse} pairs that explicitly targeted an
asset resolved through a provider, runtime binding and action, resulting
assertions, and a covering validation case; the checker reported no incomplete
chain. The two chains per asset encode the modeled answer and handoff paths,
not two observed visitor interactions.

On the natural corpus, responsibility resolution was unique at the
asset-assignment tier for 93/93 assets; none required group or zone fallback,
and none was ambiguous or unassigned. We therefore do not present fallback as
a natural-corpus frequency. In the separate derived priority suite, both
adapters produced the expected result for all nine decision cases
(9/9 per adapter, 18/18 adapter outcomes): three group fallbacks, three zone
fallbacks, and three fail-closed ambiguities.

The routing benchmark paired each asset with every possible starting agent.
The owner was already local
in 93 cases, a direct handoff was declared in 205, a route existed only
transitively in 143, and 18 were unreachable. Thus local or direct one-hop
coverage was 298/459 (64.9\%), while offline graph reachability was 441/459
(96.1\%). The latter number does not mean that the online interface performs a
multi-hop handoff in one turn. The direct-artifact and model-grounded adapters
agreed on all 93 owner decisions and all 459 route classifications under the
shared deterministic resolver.

\subsection{RQ2: Representation tradeoffs}

The natural representations had no owner or route mismatch, including after
each of nine semantically matched mutations. Appendix
Table~\ref{tab:rq2-faults} separates validation from representation edits.
Each adapter accepted all three
expected-valid case copies, so neither produced a false positive. Each detected
and localized 9/9 common mutations: missing ownership, a duplicate owner, and a
dangling handoff in every case. This result establishes parity for the injected
fault set; it does not show greater fault sensitivity for either
representation.

The edit counts expose a mixed representation tradeoff. Across the nine
patches, the model-grounded condition changed fewer top-level artifacts
(9 versus 15) but more stored references (18 versus 15). These are automatic
patch deltas, not evidence that a person would complete a maintenance task
faster or with fewer errors. For the additional V2-only mutations, the
canonical validator detected 3/9 and the stricter agent-provider contract
detected 9/9; each localized every fault it detected and yielded 0/3 false
positives on expected-valid copies. Because the direct artifacts have no
equivalent constructs, these V2-only outcomes indicate added contract
expressiveness, not a baseline failure.

Appendix Table~\ref{tab:rq2-runtime} reports local structural workload. The
fresh V2 adapter took 4.67--6.41 times the direct-adapter median. Absolute
medians remained below 8.3\,ms, but the measurement excludes parsing,
V2 generation, Unity, networking, and language-model latency. It therefore
characterizes only the in-memory adapter workload and is not an end-user
response-time result.

\subsection{RQ3: Contract enforcement and placement transactions}

The exhaustive runtime integration sweep passed all 459/459 probes. It admitted
all 298/298 local and direct routes, preserved the expected target and provider,
returned non-empty capability and runtime identifiers consistently across all
checked response fields, and made exactly 298 deterministic-stub calls. It
rejected all 161/161 transitive and unreachable routes before the stub with no
session mutation. The per-case probe totals were 135, 144, and 180; all session
states were restored and no network or external API was used. These probes
execute the request contract, but they are not 459 human interactions.

The targeted backend suite accepted the valid deictic request and rejected all
16 stale, unknown, ambiguous, contradictory, oversized, or otherwise invalid
variants without retained session mutation. Fifteen stopped before the stub;
the model-proposed unmodeled handoff required one call to reveal its destination
and was then rolled back. These are contract tests, not dialogue-quality trials.

All three Unity Editor batch entry points returned zero. Together they exercised
scenario progression, dispatch, five assertion types, valid and broken evidence,
desktop-ray selection, and ambiguous registries. The Steinpilz smoke loaded 31
contexts and 73 actions and uniquely bound 30/30 objects without UUID-suffix
fallback. These are editor batch paths, not human or headset evidence.

All four placement transactions also passed: preview did not write, undo
restored accepted bytes, stale revisions were discarded non-destructively, and
an injected regeneration fault rolled back files and session state. These
results do not imply responsibility, capability, or handoff authoring.
